% Partie : logique du premier ordre
% Chapitre : complétude
% Section : théorème de complétude

\documentclass[../../../include/open-logic-section]{subfiles}

\begin{document}

\iftag{FOL}
      {\olfileid[fr]{fol}{com}{cth}}
      {\olfileid[fr]{pl}{com}{cth}}

\olsection{Le théorème de complétude}

\begin{explain}
Rassemblons les résultats obtenus : nous arrivons au théorème de
complétude.
\end{explain}

\begin{thm}[Théorème de complétude]
\ollabel{thm:completeness}
Dans le langage au plus dénombrable fixé ici, soit $\Gamma$ un
ensemble \iftag{FOL}{d'!!{sentence}s}{de !!{formula}s}.
Si $\Gamma$ est cohérent, il est satisfaisable.
\end{thm}

\begin{proof}
Supposons $\Gamma$ cohérent.\iftag{FOL}{ D'après le
  \olref[hen]{lem:henkin}, il existe un ensemble cohérent et saturé
  $\Gamma'\supseteq\Gamma$ dans le langage enrichi~$\Lang L'$.}{} D'après le
\olref[lin]{lem:lindenbaum}, il existe un ensemble
$\Gamma^*\supseteq\iftag{FOL}{\Gamma'}{\Gamma}$ cohérent et
!!{complete}.
\iftag{FOL}{
  Puisque $\Gamma'\subseteq\Gamma^*$, pour toute
  !!{formula}~$!A(x)$ dont les variables libres appartiennent
  à~$\{x\}$, $\Gamma^*$ contient !!a{sentence} de la forme
  \iftag{prvEx}
        {$\lexists[x][!A(x)]\lif !A(c)$}
        {$\lnot\lforall[x][!A(x)]\lif\lnot !A(c)$} ;
  $\Gamma^*$ est donc saturé. Si aucun énoncé de $\Gamma$ ne
  contient~$\eq$, alors, d'après le}{D'après le}
\olref[mod]{lem:truth},
$\iftag{FOL}{\Sat{M(\Gamma^*)}}{\pSat{v(\Gamma^*)}}{!A}$ si et
seulement si $!A\in\Gamma^*$ pour \iftag{FOL}{tout !!{sentence}}{toute !!{formula}}~$!A$ de~$\Gamma$.
Tous les éléments de~$\Gamma$ sont donc satisfaits.\iftag{FOL}{
  Si un énoncé de~$\Gamma$ contient~$\eq$, alors, d'après le
  \olref[ide]{lem:truth}, pour tout !!{sentence}~$!A$,
  $\Sat{\equivclass{M}{\approx}}{!A}$ si et seulement si
  $!A\in\Gamma^*$. En particulier,
  $\Sat{\equivclass{M}{\approx}}{!A}$ pour tout énoncé
  $!A\in\Gamma$. Dans les deux cas, le réduit de la structure au
  langage initial~$\Lang L$ satisfait~$\Gamma$.}{}
Ainsi, $\Gamma$ est satisfaisable.
\end{proof}

\begin{cor}[Théorème de complétude, seconde formulation]
\ollabel{cor:completeness}
Soient $\Gamma$ un ensemble \iftag{FOL}{d'!!{sentence}s}{de !!{formula}s}
et $!A$ \iftag{FOL}{un !!{sentence}}{une !!{formula}} du même langage.
Si $\Gamma\Entails !A$, alors $\Gamma\Proves !A$.
\end{cor}

\begin{proof}
Les $\Gamma$ du~\olref{cor:completeness} et du
\olref{thm:completeness} sont quantifiés universellement. Pour éviter
toute confusion, reformulons le \olref{thm:completeness} avec une autre
variable : pour tout ensemble \iftag{FOL}{d'!!{sentence}s}{de !!{formula}s}~$\Delta$, si $\Delta$
est cohérent, alors il est satisfaisable. Par contraposition, si
$\Delta$ n'est pas satisfaisable, alors il est incohérent. Nous allons
appliquer ce résultat au corollaire.

Supposons $\Gamma\Entails !A$. Alors
$\Gamma\cup\{\lnot !A\}$ n'est pas satisfaisable, d'après la
\olref[syn][sem]{prop:entails-unsat}. En prenant
$\Delta=\Gamma\cup\{\lnot !A\}$, la première formulation du
\olref{thm:completeness} montre que cet ensemble est incohérent.
D'après les
\iftag{FOL}{%
  \tagrefs{prfAX/{fol:axd:prv:prop:prov-incons},
    prfSC/{fol:seq:prv:prop:prov-incons},
    prfND/{fol:ntd:prv:prop:prov-incons},
    prfTab/{fol:tab:prv:prop:prov-incons}}}{%
  \tagrefs{prfAX/{pl:axd:prv:prop:prov-incons},
    prfSC/{pl:seq:prv:prop:prov-incons},
    prfND/{pl:ntd:prv:prop:prov-incons},
    prfTab/{pl:tab:prv:prop:prov-incons}}},
$\Gamma\Proves !A$.
\end{proof}

\tagprob{FOL}
\begin{prob}
Utilisez le \olref[fol][com][cth]{cor:completeness} pour démontrer le
\olref[fol][com][cth]{thm:completeness}, et montrez ainsi que les deux
formulations du théorème de complétude sont équivalentes.
\end{prob}
\tagendprob

\tagprob{notFOL}
\begin{prob}
Utilisez le \olref[pl][com][cth]{cor:completeness} pour démontrer le
\olref[pl][com][cth]{thm:completeness}, et montrez ainsi que les deux
formulations du théorème de complétude sont équivalentes.
\end{prob}
\tagendprob

\tagprob{FOL}
\begin{prob}
Pour qu'un système de !!{derivation} soit complet, ses règles doivent
permettre de démontrer l'incohérence de tout ensemble non
satisfaisable. Quelles règles de !!{derivation} ont été nécessaires
pour démontrer la complétude ? Certaines de ces règles ne sont-elles
utilisées nulle part dans la démonstration ? Pour répondre, dressez
une liste ou un diagramme indiquant quelles règles de !!{derivation}
interviennent dans quels résultats menant à la démonstration du
\olref[fol][com][cth]{thm:completeness}. Relevez aussi tout emploi
implicite des règles dans ces démonstrations.
\end{prob}
\tagendprob

\tagprob{notFOL}
\begin{prob}
Pour qu'un système de !!{derivation} soit complet, ses règles doivent
permettre de démontrer l'incohérence de tout ensemble non
satisfaisable. Quelles règles de !!{derivation} ont été nécessaires
pour démontrer la complétude ? Certaines de ces règles ne sont-elles
utilisées nulle part dans la démonstration ? Pour répondre, dressez
une liste ou un diagramme indiquant quelles règles de !!{derivation}
interviennent dans quels résultats menant à la démonstration du
\olref[pl][com][cth]{thm:completeness}. Relevez aussi tout emploi
implicite des règles dans ces démonstrations.
\end{prob}
\tagendprob

\end{document}
